\chapter{Reflection, Ballot Counts and Arcsine Laws}
\label{sec:p15-reflection-and-ballot-counts}

Reflection counts paths by an explicit bijection. The ballot theorem
then counts strict leads, return-time identities count excursions, and
two different decompositions give the same arcsine distribution.
This chapter retains the combinatorial presentation of Durrett's
Section 4.9, including its renewal exercise
\cite[pp.~262--267]{DurrettPTEAuthor2019}.
It is distinct from the reverse-martingale ballot argument and the
optional-stopping generating-function calculations supplied earlier.

\begin{hbdefinition}{def:p15-rb-paths}{Paths, counts, and lattice conventions}
A length-\(n\) path from \(x\) to \(y\), where
\(n\in\mathbb N_0\) and \(x,y\in\mathbb Z\), is a sequence
\((s_0,\ldots,s_n)\) with \(s_0=x\), \(s_n=y\), and
\(s_j-s_{j-1}\in\{-1,1\}\).
Its polygonal interpolation joins consecutive vertices by straight
segments. Hitting a lattice level means visiting it at an integer time.
Here \(\mathbb N_0=\{0,1,2,\ldots\}\).
Write \(N_{n,z}\) for the number of paths from zero to \(z\).
For \(n\geq0\), set \(\binom nr=0\) unless \(r\) is an
integer with \(0\leq r\leq n\); in that range it has its
usual factorial value. In particular \(N_{0,0}=1\) and
there is exactly one empty path.

For probability assertions, \(S_0=0\) and
\(S_n=\xi_1+\cdots+\xi_n\), where the increments are
independent fair signs. Thus every length-\(n\) path has
probability \(2^{-n}\). Put
\[
 p_{n,z}=\mathbb P(S_n=z),\qquad
 u_n=\mathbb P(S_n=0).
 \]
For integers \(j\geq1\), define the level hitting time
\(T_j=\inf\{n\geq0:S_n=j\}\).
Define \(R=\inf\{n\geq1:S_n=0\}\), the first \emph{positive}
return time, and \(R_0=0\),
\[
 R_m=\inf\{n>R_{m-1}:S_n=0\}\quad(m\geq1).
 \]
An empty infimum is \(\infty\); after an infinite return time
all subsequent return times are infinite.
\end{hbdefinition}

\section{Reflection and ballot words}
\label{sec:p15-rb-reflection}

\begin{hbtheorem}{thm:p15-rb-reflection}{proved=annex:p15.rb-reflection}{The path count and the reflection bijection}
For every \(n\geq0\) and \(z\in\mathbb Z\),
\[
 N_{n,z}=\binom n{(n+z)/2},\qquad
 p_{n,z}=2^{-n}N_{n,z}.
 \]
The count is zero unless \(|z|\leq n\) and \(n-z\) is even.
For positive integers \(x,y\), reflection up to the first hit
of zero bijects the length-\(n\) paths from \(x\) to \(y\)
that hit zero with all length-\(n\) paths from \(-x\) to \(y\).
Consequently
\[
 \#\{\text{paths from }x\text{ to }y\text{ staying strictly positive}\}
                  =N_{n,y-x}-N_{n,y+x}.
 \]
This includes the zero-count and length-zero cases.
\end{hbtheorem}

\begin{hbtheorem}{thm:p15-rb-ballot}{proved=annex:p15.rb-ballot}{Strict and weak ballot inequalities}
Let \(\alpha,\beta\) be nonnegative integers,
\(n=\alpha+\beta\geq1\), and choose uniformly among all
words containing \(\alpha\) letters \(A\) and \(\beta\) letters \(B\).
A strict lead means that every nonempty prefix contains more
\(A\)'s than \(B\)'s; a weak lead permits equality.
If \(\alpha>\beta\), the strict-lead probability is
\[
                       \frac{\alpha-\beta}{\alpha+\beta}.
 \]
If \(\alpha\leq\beta\), it is zero.
If \(\alpha\geq\beta\), the weak-lead probability is
\[
                       \frac{\alpha-\beta+1}{\alpha+1};
 \]
if \(\alpha<\beta\), it is zero.
The strict formula is source Theorem 4.9.2, proved by its
reflection count; the weak formula is a separately proved boundary
comparison. The source introduces Theorem 4.9.2 as a consequence of
its reflection principle, Theorem 4.9.1, and as a result first proved
in 1878 \cite[p.~263]{DurrettPTEAuthor2019}. The source names neither
the person who proved it nor the place of publication, and neither is
supplied here. At \(\alpha=\beta=m\), the weak-lead count is
\(\binom{2m}{m}/(m+1)\). The empty word, excluded by \(n\geq1\),
satisfies both prefix conditions vacuously.
\end{hbtheorem}

\section{Returns and renewal}
\label{sec:p15-rb-returns}

\begin{hbtheorem}{thm:p15-rb-return-counts}{proved=annex:p15.rb-return-counts}{The source's no-return identity and first-return masses}
For \(n\geq0\),
\[
 \mathbb P(R>2n)=u_{2n}
                  =\binom{2n}{n}2^{-2n}.
 \]
For \(n\geq1\),
\[
 \mathbb P(S_1>0,\ldots,S_{2n}>0)=\tfrac12u_{2n},
 \qquad
 f_{2n}:=\mathbb P(R=2n)=\frac{u_{2n}}{2n-1}.
 \]
Odd first-return times have zero probability.
Moreover \(R\) has the same law as \(1+T_1\), and
\[
 \mathbb P(T_1=2n-1)=f_{2n},\qquad
 u_{2n}\sim\frac1{\sqrt{\pi n}},\qquad
 f_{2n}\sim\frac1{2\sqrt\pi\,n^{3/2}}.
 \]
Both \(R,T_1\) are finite almost surely, have infinite mean,
and, for \(p>0\), have finite \(p\)-moment exactly when \(p<1/2\).
The proof obtains these facts from the ballot count and
Stirling's formula, independently of optional stopping.
\end{hbtheorem}

\begin{hbdefinition}{def:p15-rb-renewal-chain}{The Markov premise of the source renewal exercise}
Let \((E,\mathcal E)\) be a measurable state space and fix
\(a\in E\) with \(\{a\}\in\mathcal E\).
A transition kernel \(P\) assigns a
probability measure \(P(x,\cdot)\) to each \(x\), with
\(x\mapsto P(x,B)\) measurable for every \(B\in\mathcal E\).
A time-homogeneous Markov chain \(X_n\) with this kernel satisfies
\[
 \mathbb P(X_{n+1}\in B\mid\sigma(X_0,\ldots,X_n))=P(X_n,B).
\]
The identity is required almost surely for every
\(B\in\mathcal E\) at every finite \(n\).
For a countable state space with all subsets measurable, this is
the matrix condition \(P(i,j)\geq0\), \(\sum_jP(i,j)=1\).
The law \(\mathbb P_a\) starts with \(X_0=a\).
The exercise assumes such a chain; no assertion of existence
for unspecified transition data is needed here.
Its first positive return is
\(\tau_a^+=\inf\{n\geq1:X_n=a\}\).
Put \(f_n=\mathbb P_a(\tau_a^+=n)\), \(n\geq1\),
and \(v_n=\mathbb P_a(X_n=a)\), \(n\geq0\), so \(v_0=1\).
This \(v_n\) distinguishes the general-chain probabilities from
the simple-walk \(u_n\).
\end{hbdefinition}

\begin{hbtheorem}{thm:p15-rb-renewal}{proved=annex:p15.rb-renewal}{The complete renewal equation and generating functions}
For the preceding chain and every \(n\geq1\),
\[
 v_n=\sum_{m=1}^n f_m v_{n-m}.
 \]
For \(0\leq s<1\), define
\[
 V(s)=\sum_{n\geq0}v_ns^n,\qquad
 F(s)=\sum_{n\geq1}f_ns^n.
 \]
These series converge, \(F(s)\leq s<1\), and
\[
                         V(s)=\frac1{1-F(s)}.
 \]
The same identity holds for complex \(|s|<1\), with the
series interpreted by absolute convergence.
At \(s=1\), if \(F(1)=\mathbb P_a(\tau_a^+<\infty)<1\),
then \(\sum_nv_n=1/(1-F(1))\); if \(F(1)=1\), the sum is infinite.
For the fair walk and \(0\leq s<1\), this gives
\[
 \sum_{n\geq0}u_ns^n=(1-s^2)^{-1/2},\qquad
 \mathbb E(s^R;R<\infty)=1-\sqrt{1-s^2}.
 \]
This proves both parts of Exercise 4.9.1, including the
positive-return convention and the actual regenerative premise.
\end{hbtheorem}

\section{Two different arcsine decompositions}
\label{sec:p15-rb-arcsine}

\begin{hbdefinition}{def:p15-rb-last-zero-and-occupation}{Last zero and segments on the positive side}
For integer \(n\geq0\), let
\[
 L_{2n}=\max\{0\leq j\leq2n:S_j=0\},
 \qquad
 A_{2n}=\sum_{j=1}^{2n}\mathbf1_{\{S_{j-1}+S_j>0\}}.
 \]
The first maximum is nonempty because \(S_0=0\).
An edge contributes to \(A_{2n}\) exactly when its entire
segment lies in the closed upper half-plane. In particular
the steps \(0\to1\) and \(1\to0\) contribute, while
\(0\to-1\) and \(-1\to0\) do not.
This is the source's variable \(\pi_{2n}\).
It is not the count of strictly positive integer-time vertices.
The source presents \(\pi_{2n}\) as the amount of time one player is
ahead \cite[p.~266]{DurrettPTEAuthor2019}. In the betting reading of
\cref{thm:p15-rb-calculations}, item 4, where \(S_j\) is the first
player's net gain after \(j\) unit bets, \(A_{2n}\) is the number of
the \(2n\) betting periods during which that player is ahead in the
polygonal sense: at the start and at the end of the period the gain is
nonnegative, and it is positive at one of the two.
Both \(L_{2n}\) and \(A_{2n}\) take values in
\(\{0,2,\ldots,2n\}\); the proof establishes the occupation parity.
\end{hbdefinition}

\begin{hbtheorem}{thm:p15-rb-last-zero}{proved=annex:p15.rb-last-zero}{The last-zero decomposition}
For \(0\leq k\leq n\),
\[
                   \mathbb P(L_{2n}=2k)=u_{2k}u_{2n-2k}.
 \]
In particular the finite distribution is symmetric in
\(k\) and \(n-k\), and
\(\sum_{k=0}^n u_{2k}u_{2n-2k}=1\).
The endpoint values \(k=0,n\) use \(u_0=1\).
\end{hbtheorem}

\begin{hbtheorem}{thm:p15-rb-arcsine-limit}{proved=annex:p15.rb-arcsine-limit}{The last-zero arcsine law with endpoint control}
As \(n\to\infty\), \(L_{2n}/(2n)\) converges weakly to the
probability law on \([0,1]\) with density
\[
 h(x)=\frac1{\pi\sqrt{x(1-x)}}\quad(0<x<1),
 \]
with \(h(x)=0\) outside \((0,1)\) and no endpoint atoms.
Its CDF is \(H(x)=2\pi^{-1}\arcsin\sqrt x\) for
\(0\leq x\leq1\), zero below zero and one above one.
For every \(0<a<b<1\),
\[
 \mathbb P(a\leq L_{2n}/(2n)\leq b)
                                  \longrightarrow\int_a^b h(x)\,dx.
 \]
The proof retains the source's step-density/Riemann-integral route,
including the exact rounded indices, and then proves the endpoint
extension. It does not assume a uniform approximation at the
singular endpoints.
\end{hbtheorem}

\begin{hbtheorem}{thm:p15-rb-occupation}{proved=annex:p15.rb-occupation}{The separate first-excursion occupation proof}
For every \(0\leq k\leq n\),
\[
                  \mathbb P(A_{2n}=2k)=u_{2k}u_{2n-2k}.
 \]
Consequently \(A_{2n}/(2n)\) has the same arcsine weak limit
as \(L_{2n}/(2n)\). The proof uses the source's induction on
the first complete excursion, with the two endpoint cases
proved separately. Equality of these two laws does not assert
pathwise equality of the variables.
\end{hbtheorem}

\section{The return-time limit stated in the source}
\label{sec:p15-rb-half-stable}

\begin{hbtheorem}{thm:p15-rb-laplace-premise}{proved=annex:p15.rb-laplace-premise}{A positive half-stable law and a sufficient transform criterion}
The function
\[
 g(y)=
 \begin{cases}
  (2\pi)^{-1/2}y^{-3/2}e^{-1/(2y)},&y>0,\\
  0,&y\leq0
 \end{cases}
 \]
is a probability density. If \(Y\) has this density, then
\[
                    \mathbb Ee^{-\lambda Y}=e^{-\sqrt{2\lambda}}
                                      \quad(\lambda\geq0).
 \]
Laplace transforms on \([0,\infty)\) determine probability laws.
If finite nonnegative \(Z_n\) satisfy
\(\mathbb Ee^{-\lambda Z_n}\to\mathbb Ee^{-\lambda Y}\)
for every \(\lambda\geq0\), then \(Z_n\Rightarrow Y\).
Moreover \(Y_1+\cdots+Y_m\) has the same law as \(m^2Y\)
for independent copies \(Y_j\); this defines its positive
strictly stable law of index \(1/2\).
The proof constructs and normalizes the density and proves
every transform assertion.
\end{hbtheorem}

\begin{hbtheorem}{thm:p15-rb-return-scaling}{proved=annex:p15.rb-return-scaling}{Returns and level hits on their exact square scale}
All \(R_m,T_m\), \(m\geq1\), are finite almost surely, and
\[
                  R_m\ \stackrel{\mathrm d}=\,m+T_m.
 \]
As \(m\to\infty\),
\[
                      R_m/m^2\Rightarrow Y,\qquad
                      T_m/m^2\Rightarrow Y,
 \]
where \(Y\) has the preceding density.
The finite-time return and level-hit laws are different; they
have the same limit on this scale.
This supplies the exact return-time limit asserted after
source equation (4.9.2), using the combinatorially derived
return generating function.
\end{hbtheorem}

\begin{hbtheorem}{thm:p15-w02c-p15-020-return-domain}{proved=annex:p15-w02c-p15-020-return-domain}{The return time lies in the domain of attraction of the index one-half law}
Say that a random variable \(X\) lies in the \emph{domain of
attraction} of a nondegenerate law \(Y\) if, for independent copies
\(X_1,X_2,\ldots\) of \(X\), there are constants \(a_m>0\) and
\(b_m\) with \((X_1+\cdots+X_m-b_m)/a_m\Rightarrow Y\); this is the
convergence in the source's Theorem~3.8.2 and its remark
\cite[Section~3.8, p.~186]{DurrettPTEAuthor2019}. Let \(Y\) be as in
\cref{thm:p15-rb-laplace-premise}.
\begin{enumerate}
\item \(R\) and \(T_1\) lie in the domain of attraction of \(Y\), with
 \(a_m=m^2\) and \(b_m=0\) for \(R\), and \(a_m=m^2\), \(b_m=0\) for
 \(T_1\) as well.
\item \(Y\) is stable in the source's sense: for each integer
 \(k\geq1\), independent copies satisfy
 \((Y_1+\cdots+Y_k)/k^2\stackrel{\mathrm d}=Y\)
 \cite[Section~3.8, p.~191]{DurrettPTEAuthor2019}. Its characteristic
 function is
 \[
  \mathbb Ee^{itY}=\exp\{-|t|^{1/2}(1-i\operatorname{sgn}t)\}
  \qquad(t\in\mathbb R),
 \]
 so the scale \(k^2=k^{1/\alpha}\) and the exponent \(|t|^{1/2}\) both
 give the index \(\alpha=1/2\).
\item The source's skewness parameter is \(\kappa=2\theta-1\), where
 \(\theta=\lim_{x\to\infty}\mathbb P(X>x)/\mathbb P(|X|>x)\)
 \cite[Section~3.8, pp.~186 and 189]{DurrettPTEAuthor2019}. Since
 \(R\geq1\), \(\theta=1\) for \(R\), so \(\kappa=1\): all of the tail
 lies on the positive side, and \(Y>0\) almost surely.
\end{enumerate}
This is the source's conclusion that \(R\) lies in the domain of
attraction of the stable law with \(\alpha=1/2\) and \(\kappa=1\),
which it draws from (4.9.2) \cite[p.~264]{DurrettPTEAuthor2019}.
The source also records that the law with \(\alpha=1/2\),
\(\kappa=1\), \(c=0\) and \(b=1\) has the density \(g\) of
\cref{thm:p15-rb-laplace-premise}, its display (3.8.12)
\cite[Section~3.8, p.~189]{DurrettPTEAuthor2019}. Matched literally
against the source's printed form (3.8.11) of the characteristic
function, the function in item 2 has \(c=0\), \(b=1\) and
\(w_{1/2}=\tan(\pi/4)=1\), but its imaginary part carries the sign that
the printed form gives to \(\kappa=-1\). The value \(\kappa=1\) used
here is therefore the source's tail value \(2\theta-1\) together with
its density (3.8.12), not a literal reading of (3.8.11).
\end{hbtheorem}

\section{Worked counts and interpretation boundaries}
\label{sec:p15-rb-calculations}

\begin{hbtheorem}{thm:p15-rb-calculations}{proved=annex:p15.rb-calculations}{Explicit counts, finite moments and comparisons}
\begin{enumerate}
\item There are six length-four paths from \(1\) to \(1\);
 four hit zero and two remain positive.
 With three \(A\) votes and two \(B\) votes, exactly two of
 the ten words give a strict lead and five give a weak lead.
\item For four fair steps, both \(L_4\) and \(A_4\) have
 probabilities \(3/8,1/4,3/8\) at \(0,2,4\).
 The path \((0,1,0,-1,0)\) has \(L_4=4\) and \(A_4=2\).
 The path \((0,1,0,1,0)\) has four positive-side edges but
 only two strictly positive sampled vertices.
\item For \(n\geq1\), either \(K=L_{2n}/2\) or \(K=A_{2n}/2\)
 has
 \[
  \mathbb EK=n/2,\qquad \operatorname{Var}(K)=n(n+1)/8.
 \]
 Thus either normalized fraction has mean \(1/2\) and
 variance \((n+1)/(8n)\to1/8\).
 The largest finite masses are at the endpoints, and the
 smallest are at the central one or two admissible indices.
\item In the limiting arcsine law, the middle interval
 \([1/4,3/4]\) has probability \(1/3\), while the two outer
 quarters together have probability \(2/3\).
 The probability of no zero after the halfway time tends to
 \(1/2\). The source gives this a betting reading: two players
 stake one unit on a fair coin flip every day, and \(S_j\) is the
 first player's net gain after \(j\) days
 \cite[p.~265]{DurrettPTEAuthor2019}. Over \(2n\) days, the event
 \(L_{2n}\leq n\) is exactly the event that one of the two players
 is strictly ahead at the end of every day from day \(n+1\) to day
 \(2n\), the same player throughout. So the probability that one
 player leads for the whole second half tends to \(1/2\); the source
 phrases this for a year, with the lead held from July 1 to the end of
 the year. These are asymptotic statements, not exact calendar-year
 probabilities: a year has an odd number of days, and the limit
 concerns \(n\to\infty\).
\item The limiting arcsine density \(h\) attains its smallest value,
 \(2/\pi\), at \(x=1/2\) and only there, it decreases on
 \((0,1/2]\) and increases on \([1/2,1)\), and \(h(x)\to\infty\) as
 \(x\to0\) or \(x\to1\). Read through \(A_{2n}\), an equal division
 of the betting periods between the two players is the least likely
 region of the limit law, and nearly one-sided divisions are the most
 likely; this is the source's remark after its Theorem 4.9.6
 \cite[p.~266]{DurrettPTEAuthor2019}. The finite-\(n\) counterpart is
 the mass profile in item 3.
\item Reflection preserves path counts but need not preserve
 biased path probabilities. Let \(0<p<1\) be the probability
 of an up-step and \(q=1-p\). For two steps from \(1\) to \(1\),
 the touching path has probability \(qp\), whereas its reflected
 path from \(-1\) to \(1\) has probability \(p^2\).
\item The renewal identity needs its stated homogeneous Markov
 premise and a positive return time. An arbitrary deterministic
 process need not satisfy it; first hitting including time zero
 is not a valid replacement for the return time.
\end{enumerate}
\end{hbtheorem}

\begin{proofboundary}{def:p15-rb-next-boundary}{Completed combinatorics and the remaining source routes}
The numbered content of Section 4.9, its sole Exercise 4.9.1,
and its stated return-time limit are supplied.
Here that limit is proved by the direct transform calculation of
\cref{thm:p15-rb-return-scaling}. The source reaches the same
limit by two further routes.

The first route places the return time in the domain of
attraction of the stable law with index \(1/2\), from the tail
asymptotics (4.9.2) and the source's Theorem~3.8.8
\cite[p.~264]{DurrettPTEAuthor2019}; that theorem, which
characterizes the possible limits of normalized i.i.d.\ sums as
the stable laws, is stated and proved at
\cite[Section~3.8, p.~191]{DurrettPTEAuthor2019}.
The conclusion of this route, with its parameters \(\alpha=1/2\)
and \(\kappa=1\), is proved in
\cref{thm:p15-w02c-p15-020-return-domain} from the exact transform
limit, without the tail route. The handbook contains no general
theory of stable laws or of their domains of attraction, so the
route itself, from tail asymptotics to a stable limit, is not
reproduced. Its only general domain-of-attraction chapter,
\cref{sec:p15-normal-domains-of-attraction}, concerns the normal
law alone, and its criterion does not apply to a limit of index
\(1/2\).

The second route identifies the limit with the first hitting
time of level \(1\) by Brownian motion. The source announces it
after (4.9.2) \cite[pp.~264--265]{DurrettPTEAuthor2019} and
carries it out in Example~8.1.12, through its invariance
principle for random walk paths
\cite[Section~8.1, p.~395]{DurrettPTEAuthor2019}. The handbook
constructs neither Brownian motion nor an invariance principle,
so this identification is not proved here.
The density \(g\) of \cref{thm:p15-rb-laplace-premise} coincides
with the case \(a=1\) of the hitting-time density in the source's
display (7.4.6) \cite[Section~7.4, p.~374]{DurrettPTEAuthor2019}.
That coincidence is a comparison of two formulas; nothing about
Brownian motion is derived from it, and the half-stable
calculation above does not depend on it.
\end{proofboundary}
